\documentclass[12pt]{article}
\usepackage{enumitem}
\usepackage{amsmath}
\usepackage{fullpage}
\usepackage[T1]{fontenc}   % better ASCII glyphs
\usepackage{inconsolata}   % nicer monospaced font (zi4)

\begin{document}

\begin{center}
\LARGE\textbf{Practice Exam Questions}
\end{center}

\normalsize

\begin{enumerate}[label=\textbf{\arabic*.}]

\item Write a clear one--sentence definition of Level I, Level II, and Level III regression as described by Berk (2010).
\newpage
% Solution: Level I: descriptive, summarize associations; 
% Level II: inference to population with SEs/CIs, no causality; 
% Level III: causal inference with credible identification.

\item Suppose you have a fitted logit model. Define a first difference using both words and equations. Be precise! Describe how to interpret the first difference.

% Solution: FD = E(y | X_hi) - E(y | X_lo) = Pr(y=1 | X_hi) - Pr(y=1 | X_lo).

\newpage

\item In the Clark and Golder (2006) normal model exercise, you fit the normal model by maximum likelihood using \verb|optim()|. The parameter vector \(\theta\) stacks the \(k\) regression coefficients \(\beta\) (for the \(k\) columns in the design matrix \(X\)) and then the scalar \(\sigma\). Which extraction is correct?
\begin{itemize}
  \item[A.] \verb|beta <- theta[1:(k-1)]; sigma <- theta[k]|
  \item[B.] \verb|beta <- theta[1:k]; sigma <- theta[k+1]| % correct
  \item[C.] \verb|beta <- theta[2:(k+1)]; sigma <- theta[1]|
  \item[D.] \verb|beta <- theta; sigma <- 1|
\end{itemize}

\item \textit{Select all that apply.} Which of the following formulas \textbf{exclude} or \textbf{omit} an intercept term from the design matrix?
\begin{itemize}
  \item[A.] \verb|~ 1 + x1 + x2| % correct
  \item[B.] \verb|~ -1 + x1 + x2|
  \item[C.] \verb|~ 0 + x1 + x2|
  \item[D.] \verb|~ x1 + x2| % correct
\end{itemize}

\item Suppose \verb|species| is a factor with 3 categories. By default, \verb|model.matrix(~ species)| creates:
\begin{itemize}
  \item[A.] Three indicator columns, one for each category
  \item[B.] Two indicator columns (treatment coding), with one category as the baseline % correct
  \item[C.] No columns unless you manually expand
  \item[D.] A single numeric column with levels 1, 2, 3
\end{itemize}

\item Including \verb|poly(x, 2)| in a formula adds columns \(x\) and \(x^2\) to the design matrix.
\begin{itemize}
  \item[A.] TRUE
  \item[B.] FALSE
\end{itemize}

\item Consider the formula \verb|y ~ x * z|. What predictors does this generate in the design matrix?
\begin{itemize}
  \item[A.] Only the product term \verb|x| $\times$ \verb|z|
  \item[B.] \verb|x|, \verb|z|, and their product \verb|x| $\times$ \verb|z|% correct
\end{itemize}

\item Suppose \verb|f| is a formula. Which pipeline ensures that \(X\) and \(y\) are properly constructed on the same rows (e.g., after dropping missing values from variables used in the model)?
\begin{itemize}

  \item[A.]
  \begin{verbatim}
X <- model.matrix(f, turnout)
y <- turnout$vote
  \end{verbatim}

  \item[B.]
  \begin{verbatim}
mf <- model.frame(f, data = turnout)
X  <- model.matrix(f, mf)
y  <- model.response(mf)
  \end{verbatim} % correct

  \item[C.]
  \begin{verbatim}
X <- model.matrix(f, na.omit(turnout))
y <- na.omit(turnout)$vote
  \end{verbatim}

\end{itemize}


\item When fitting the logit model via \verb|optim()| with design matrix \(X\), what is a correct choice for the starting parameter vector?
\begin{itemize}
  \item[A.] \verb|par_start <- rep(0, nrow(X))|
  \item[B.] \verb|par_start <- rep(0, ncol(X))| % correct
  \item[C.] \verb|par_start <- 0|
  \item[D.] \verb|par_start <- rnorm(nrow(X))|
\end{itemize}

\item In the probit--logit comparison exercise, you learned the ``\(\approx 1.6\) rule-of-thumb.'' The \underline{\hspace{2cm}} are about 1.6 times larger in
\underline{\hspace{2cm}} models compared to \underline{\hspace{2cm}} models.

\begin{itemize}
  \item[A.] coefficients; logit; probit % correct
  \item[B.] coefficients; probit; logit
  \item[C.] first differences; logit; probit
  \item[D.] first differences; probit; logit
\end{itemize}

\item \textit{Choose all that apply.} In the Clark and Golder (2006) \textit{t}-model, you estimated \((\beta,\sigma,\nu)\) with ML using a location--scale \textit{t} distribution. Which statements are true?
\begin{itemize}
  \item[A.] With small \(\nu\) the \textit{t} has heavier tails than the normal, reducing outlier influence. % correct
  \item[B.] As \(\nu \to \infty\), the \textit{t} approaches the normal as a limiting case. % correct
  \item[C.] In a location--scale \textit{t} regression, the link for \(\mu = X\beta\) is still the identity; no special inverse link is introduced. % correct
  \item[D.] The \textit{t} model introduces an additional parameter \(\nu\); the normal model is effectively the \textit{t} model with \(\nu = \infty\). % correct
  \item[E.] The \textit{t} model always yields identical \(\hat{\beta}\) to OLS.
  \item[F.] Smaller \(\nu\) implies lighter tails than the normal.
\end{itemize}

\end{enumerate}

\end{document}
